\begin{helper}
Fix a set \(I\) and a sample \(\omega\).  In the red case, for every
\(r,s\in[m]\) and every ordered pair \(q\), if
\(q=(r,s)\) or \(q=(s,r)\), if \(q_1<q_2\), if
\(\noderef{TwoBiteStageF1}(\omega,I,\operatorname{red}(s))\), and if
\(\noderef{TwoBiteProjectionContains}(\omega,I,\operatorname{red}(r))\), then
\(\operatorname{red}(q)\in
\noderef{TwoBitePreTerminalRecordedPairs}(\omega,\varepsilon,I)\).  The blue
case is identical: for every \(b,c\in[m]\) and ordered pair \(q\), if
\(q=(b,c)\) or \(q=(c,b)\), if \(q_1<q_2\), if
\(\noderef{TwoBiteStageF1}(\omega,I,\operatorname{blue}(c))\), and if
\(\noderef{TwoBiteProjectionContains}(\omega,I,\operatorname{blue}(b))\), then
\(\operatorname{blue}(q)\in
\noderef{TwoBitePreTerminalRecordedPairs}(\omega,\varepsilon,I)\).
\end{helper}
\begin{proof}
Consider the red case.  Since \(s\in F_1\),
\(\noderef{TwoBiteStageRevealedVertex}(\omega,\varepsilon,I,\operatorname{red}(s))\)
holds by the \(F_1\)-branch in the definition of staged-revealed vertices.
Together with projection containment of \(\operatorname{red}(r)\), this is
exactly the incident-coordinate branch in the definition of
\(\noderef{TwoBitePreTerminalRecordedPairs}\), after putting the unordered
pair \(\{r,s\}\) into increasing order \(q\).  The hypothesis
\(q_1<q_2\) supplies membership of the tagged coordinate in
\(\noderef{TwoBiteTerminalCoordinateUniverse}(m)\), so the outer terminal
filter is also satisfied.

The blue case is word-for-word the same, using the blue \(F_1\)-branch and
the blue incident-coordinate branch of
\(\noderef{TwoBitePreTerminalRecordedPairs}\).
\end{proof}
